\section{Related work}
\paragraph{Learning with incomplete modalities.}
Representation, alignment and fusion are distinct challenges in multimodal learning \citep{baltrusaitis2019multimodal}. Modality dropout suppresses modalities during training so that a model learns to tolerate absent inputs \citep{neverova2016moddrop}. TIMM introduces type-induced missing-aware mixture-of-experts ideas \citep{qian2025timm}; Flex-MoE uses a missing-modality bank with generalized and specialized expert learning \citep{yun2024flexmoe}; UniMIX-AD develops graph and expert mechanisms for incomplete and imbalanced Alzheimer's diagnosis \citep{yang2026unimix}. We adapt mechanisms from these approaches to a common regional-input protocol. Their original encoders, losses and routing curricula are not fully reproduced, so the reported scores belong to the local adaptations.

\paragraph{Imbalanced objectives.}
Class imbalance affects both model fitting and the choice of evaluation metrics \citep{he2009imbalanced}. Class-balanced losses adjust the influence of unequal label frequencies \citep{cui2019classbalanced}. Logit adjustment changes decision scores or training objectives using label priors \citep{menon2021longtail}. Group-robust learning instead focuses on performance across predefined groups \citep{sagawa2020distributionally}. Our joint cell is defined by diagnosis and observed modality pattern, and our chosen weighting rule is a smoothed inverse square root of its expected post-dropout count. This differs from the effective-number formula of \citet{cui2019classbalanced}. It also differs from optimizing the worst-group risk: weighting encourages broader coverage but does not guarantee a minimum recall in every cell.

\paragraph{Trees and fusion.}
Tree-based models remain competitive with deep networks on medium-sized tabular benchmarks \citep{grinsztajn2022trees}; LightGBM provides an efficient boosting implementation \citep{ke2017lightgbm}. A tree and a neural model may make complementary errors on the same inputs, in which case averaging their predictions can help. We retain a seed-ensembled tree reference and test a scalar fusion with alternative neural members. This comparison asks whether balancing produces errors that are useful in fusion, without assuming that the neural model must outperform the tree on its own.
